\chapter{Umbral operativo}

\section{Noción de umbral}
\begin{definition}
Se llama umbral operativo al punto mínimo a partir del cual una formulación, configuración o trayectoria puede ser tratada con una capacidad operativa más fuerte que la mera legibilidad basal.
\end{definition}

\begin{observation}
El umbral operativo no sustituye la base. Su función es señalar cuándo una estructura ya puede entrar en tratamiento comparativo, jerárquico o protocolizable sin perder dependencia canónica.
\end{observation}

\section{Condición mínima de umbral}
\begin{definition}
Una pieza alcanza umbral operativo cuando concurren simultáneamente:
\begin{enumerate}
  \item legibilidad suficiente de sus componentes;
  \item coherencia local no colapsada;
  \item acumulación mínima registrable;
  \item persistencia local suficiente para tratamiento no instantáneo;
  \item ganancia operativa discernible respecto del nivel anterior.
\end{enumerate}
\end{definition}

\begin{observation}
La concurrencia de estas condiciones no implica todavía estabilidad madura. Solo fija una primera frontera de operabilidad disciplinada.
\end{observation}

\section{Umbral y transición}
\begin{definition}
Toda transición de nivel operativamente válida debe atravesar un umbral que permita justificar el fortalecimiento de la formulación sin ruptura de herencia estructural.
\end{definition}

\begin{observation}
Sin umbral operativo explícito, una transición corre el riesgo de quedar como incremento verbal sin soporte estructural suficiente.
\end{observation}

\section{Umbral y comparación}
\begin{definition}
Dos formulaciones son comparables en sentido operativo solo si ambas han alcanzado un umbral compatible dentro del dominio declarado.
\end{definition}

\begin{observation}
Esta cláusula evita comparar de manera simétrica piezas que todavía pertenecen a niveles de elaboración no homogéneos.
\end{observation}

\section{Condición mínima de fallo}
\begin{definition}
Se considera que una pieza pierde umbral operativo cuando deja de sostener alguna de las condiciones mínimas requeridas para su tratamiento disciplinado en el nivel correspondiente.
\end{definition}

\begin{observation}
La pérdida de umbral no equivale necesariamente a destrucción total de la estructura. Puede significar solo retorno a un nivel anterior de tratamiento.
\end{observation}

\section{Cláusula de alcance}
\begin{observation}
El umbral operativo fija una frontera mínima de entrada al trabajo canónico intermedio. No constituye todavía una teoría completa de umbrales, escalas o jerarquías exhaustivas del sistema.
\end{observation}

